\documentclass[11pt,letterpaper]{article}

% PAGE FORMAT
\usepackage[letterpaper,top=3cm,bottom=3cm,left=2cm,right=2cm]{geometry}

% ENCODING AND TYPOGRAPHY
\usepackage[T1]{fontenc}
\usepackage[utf8]{inputenc}
\usepackage[spanish,es-tabla,es-nodecimaldot,es-noshorthands]{babel}
\usepackage{parskip}

% MATHEMATICS
\usepackage{amsmath,amsfonts,amssymb,amsthm}

% TABLES AND FIGURES
\usepackage{array,booktabs,multirow}
\usepackage[table]{xcolor}
\usepackage{graphicx,caption,float}

% BIBLIOGRAPHY (IEEE, escrita al final del documento)
\usepackage{cite}
\usepackage{hyperref}

\newcommand{\um}{\,\mu\mathrm{m}}
\newcommand{\uN}{\,\mu\mathrm{N}}

\title{\textbf{\LARGE Influencia de la Configuración Geométrica
de una Telaraña Orbicular en su Dinámica Vibratoria
Mediante un Modelo de Mecánica Analítica}}
\author{Amelia Hoyos Vélez, Nicolle Rodríguez Estupiñán, Miguel Ángel León Botero\\[0.3cm]
\small Ingeniería Física, Escuela de Ciencias e Ingeniería, Universidad EAFIT}
\date{Semana 12 --- 2026}

\begin{document}
\maketitle

\section{Definición del Problema}
Cuando una presa cae sobre una telaraña orbicular, impone una perturbación mecánica localizada sobre uno de sus nodos. Esa perturbación no se queda en el punto de impacto: se propaga a través de los hilos en forma de ondas que recorren toda la estructura y eventualmente llegan a las patas de la
araña. Lo que la araña es capaz de inferir ---posición, masa y movimiento de la presa--- depende directamente de cómo la red transforma esa señal mientras viaja: qué frecuencias amplifica, cuáles atenúa, y en qué dirección predomina el movimiento cuando la onda llega a cada nodo
\cite{mortimer2016,masters1981,mortimer2019decoding}.

Esta propagación está gobernada por la mecánica de la estructura. Una telaraña orbicular es, en esencia, un sistema de masas interconectadas por elementos elásticos pretensados, dispuestos sobre una geometría radial en
el espacio tridimensional. Cada hilo tiene dos roles simultáneos: transmite fuerzas de restitución elásticas proporcionales a su deformación axial, y
ejerce fuerzas de restitución geométricas ---proporcionales a su tensión inicial--- que se oponen a cualquier movimiento transversal al plano de la red. Esta segunda contribución es la responsable de las ondas fuera del plano, que son precisamente las que la araña detecta con mayor sensibilidad \cite{mortimer2019,nakata2010,lott2022,soler2016}. El objetivo del modelo matemático que se desarrolla a continuación es capturar ambos mecanismos de manera rigurosa, para después estudiar cómo cambia el comportamiento vibratorio al variar la geometría de la red.

\end{document}